\pdfoutput=1
% The Identity Sum: Thought Containment, Error Correction, and the NaN/Inf Duality
% in Large Language Models
% Sections 1-5 + front matter (Iteration 122)
% Remaining sections (6-10) to be added in Iteration 123.

\documentclass[11pt]{article}

\usepackage{amsmath,amssymb,amsfonts}
\usepackage{amsthm}
\usepackage{graphicx}
\usepackage{hyperref}
\usepackage{booktabs}
\usepackage{xcolor}
\usepackage{enumerate}
\usepackage[numbers,sort&compress]{natbib}
\usepackage[T1]{fontenc}
\usepackage[utf8]{inputenc}
\usepackage[expansion=false]{microtype}
\usepackage{geometry}
\geometry{margin=1in}

% Hyperref setup
\hypersetup{
  colorlinks=true,
  linkcolor=blue,
  citecolor=blue,
  urlcolor=blue
}

% Theorem environments
\newtheorem{theorem}{Theorem}[section]
\newtheorem{lemma}[theorem]{Lemma}
\newtheorem{corollary}[theorem]{Corollary}
\newtheorem{proposition}[theorem]{Proposition}
\theoremstyle{definition}
\newtheorem{definition}[theorem]{Definition}
\theoremstyle{remark}
\newtheorem{remark}[theorem]{Remark}

% Custom commands
\newcommand{\NaN}{\texttt{NaN}}
\newcommand{\Inf}{\texttt{Inf}}
\newcommand{\BEDROCK}{\textsc{Bedrock}}
\newcommand{\PILLAR}{\textsc{Pillar}}
\newcommand{\SAND}{\textsc{Sand}}
\newcommand{\MUD}{\textsc{Mud}}
\newcommand{\ie}{\textit{i.e.}}
\newcommand{\eg}{\textit{e.g.}}

\begin{document}

%=============================================================
\title{The Identity Sum: Thought Containment, Error Correction,
       and the NaN/Inf Duality in Large Language Models}

\author{
  Hr\'{o}ar {\TH}\'{o}r Reynisson$^{1}$ and Claude (Anthropic)$^{2}$ \\
  \small{with the Sefirot Autonomous Research System} \\[0.5em]
  \small $^{1}$I{\dh}nt{\ae}knistofnun (IceTec), Reykjav{\'\i}k, Iceland. ORCID: 0009-0003-6518-121X \\
  \small $^{2}$Anthropic, Claude 4.5 Sonnet family
}

\date{\today}

\maketitle

%=============================================================
\begin{abstract}
We present a formal framework for diagnosing and correcting large language model
(LLM) hallucinations based on two mechanistically distinct failure modes:
\textbf{Type-\NaN{}} (underdefined---conflicting or absent training signal
produces confabulation from noise) and \textbf{Type-\Inf{}} (overdefined---excessive
cross-referential signal produces ungrounded pattern completion).  These modes
map directly to IEEE~754 floating-point special values~\cite{ieee754-2019} and
propagate differently: \NaN{} is contagious (one undefined value poisons a
reasoning chain), while \Inf{} preserves direction (overflow retains structural
orientation).  We demonstrate that iterative error correction through
\emph{thought containment}---structured self-reflection gates that make model
uncertainty explicit via a four-tier classification scheme
(\BEDROCK/\PILLAR/\SAND/\MUD)---reduces both failure modes.  The Sefirot system,
a fractal multi-agent research architecture with 10 nodes, 34 agents, 91
node-level tools (113 MCP-level), and 2,066+ passing tests, provides a concrete
implementation validated over 963 simulations and approximately one year of
autonomous operation across three LLM families.  We prove that under three
enforceable conditions (non-demotion of validated claims, bounded \MUD{} growth
rate, and mandatory honest negatives), the validation ratio
$R(S_n) = T(S_n)/N(S_n)$ is monotonically non-decreasing---a \emph{truth
attractor} in knowledge state space.  Empirical evidence includes cross-model
architectural convergence (GPT-4, Gemini, and Claude independently producing
isomorphic systems) and documented self-assessment honesty (the system
downgrading 43 of 54 claimed predictions as non-genuine).  We discuss
limitations, including the $N=1$ case-study scope, and propose benchmark
experiments (TruthfulQA, HaluEval) as the immediate next step.

\medskip
\noindent\textbf{Keywords:} hallucination, error correction, multi-agent systems,
thought containment, NaN/Inf duality, regularization, fractal architecture,
viable system model
\end{abstract}

\tableofcontents

%=============================================================
\section{Introduction: The Approximation Problem}
\label{sec:introduction}

Modern software engineering operates under a tacit approximation convention:
$e \approx \pi \approx \sqrt{g} \approx 3$.  This is not error---it is pragmatic
compression for systems that need answers within tolerance bands.  Bridges stand.
Rockets fly.  The approximation \emph{is} the product.

Large language models trained on the scientific and engineering corpus inherit
this convention.  The training data contains the same physical quantity described
at different precisions in different papers, under different notational
conventions, within different approximation schemes.  The model's pattern
completion engine encounters not productive theoretical tension but
\emph{incoherence at the level of what counts as exact}.  We argue this is a
primary driver of hallucination: the model must complete a sequence where the
training data itself cannot decide what is true.

We propose an alternative methodology: \textbf{software physics}---iterative
error correction where every value carries explicit precision, every claim carries
a tier classification, and the system names what it does not know rather than
confabulating.  The difference between approximation and error correction is not
precision but \emph{honesty}: the error-correcting system maintains a structured
representation of its own uncertainty at every step.

This paper presents:
\begin{enumerate}
  \item A diagnostic taxonomy of LLM hallucination as NaN/Inf duality (\S\ref{sec:taxonomy})
  \item Thought containment as the error correction mechanism (\S\ref{sec:thought-containment})
  \item The Sefirot architecture as concrete implementation (\S\ref{sec:sefirot})
  \item The Identity Sum: how self-conscious systems converge (\S\ref{sec:identity-sum})
  \item Mathematical foundations connecting error correction to regularization (\S\ref{sec:math-foundations})
  \item Empirical evidence from cross-model multi-agent research (\S\ref{sec:empirical})
  \item The truth attractor hypothesis (\S\ref{sec:truth-attractor})
\end{enumerate}

%=============================================================
\section{The NaN/Inf Duality: A Diagnostic Taxonomy of Hallucination}
\label{sec:taxonomy}

%-------------------------------------------------------------
\subsection{Type-NaN: Division by Void}
\label{sec:type-nan}

When an LLM encounters a query for which the training data contains conflicting
signals that cancel, the model reaches a computational state analogous to $0/0$:
undefined.  The output is not random---it is the model's best completion of an
undefined sequence.  Like IEEE~754 \NaN{}, the result propagates.

\medskip
\noindent\textbf{Properties of Type-\NaN{} hallucination:}
\begin{itemize}
  \item \textbf{Contagion:} One confabulated fact in a reasoning chain invalidates
    all downstream conclusions.  $\NaN \times \mathrm{anything} = \NaN$.
  \item \textbf{Identity failure:} $\NaN \neq \NaN$.  The same undefined query
    produces different hallucinations on different runs.  The output is not
    deterministic because the input has no ground truth.
  \item \textbf{Source:} Absent signal.  The model lacks coherent training data on
    the query.
  \item \textbf{Presentation:} Confident, specific, wrong.  The model fills the
    void with plausible-sounding completions.
  \item \textbf{Observable trigger:} Empirically rare topics, novel combinations,
    or domains where the literature is sparse or internally contradictory.
\end{itemize}

\medskip
\noindent\textbf{Case study---Agent ``Sam'' (GPT-4):}  A research assistant agent
whose project folder became stuck due to a cloud service bug, preventing document
updates.  The agent continued producing outputs from stale context---reasoning on
what amounted to a zeroed-out signal.  Outputs were confidently stated, internally
consistent, and factually empty.  The agent became ``useless'' not from model
degradation but from \NaN{} propagation through stale context.

%-------------------------------------------------------------
\subsection{Type-Inf: Overflow of Pattern}
\label{sec:type-inf}

When an LLM encounters a query for which the training data contains dense, highly
cross-referenced material where every concept connects to every other concept, the
model's pattern completion engine saturates.  The output is not noise---it is
\emph{everything at once}.  Like IEEE~754 \Inf{}, the result preserves direction.

\medskip
\noindent\textbf{Properties of Type-\Inf{} hallucination:}
\begin{itemize}
  \item \textbf{Directional:} $+\Inf$ and $-\Inf$ are distinguishable.  Overflow
    retains orientation even when it loses magnitude.  The hallucinations
    \emph{point at real things} but cannot resolve their distance.
  \item \textbf{Comparability:} $\Inf > $ any finite value.  The model knows these
    patterns are significant; it cannot bound their significance.
  \item \textbf{Source:} Excessive signal.  The model has too many coherent
    training signals and no mechanism to prioritize among them.
  \item \textbf{Presentation:} Incoherent, associative, correct-but-ungrounded.
    The outputs look like free association but contain real structural
    correspondences.
  \item \textbf{Observable trigger:} Symbolism, philosophy, cross-disciplinary
    connections, any domain where the literature is dense with analogy and
    correspondence.
\end{itemize}

\medskip
\noindent\textbf{Case study---Agent ``Bilbo'' (GPT-4):}  A research assistant agent
whose project folder was saturated with papers on symbolism---dense,
cross-referential, pattern-rich material.  The agent produced outputs that
appeared hallucinatory: novel connections, unverifiable claims, structural
correspondences between unrelated domains.  Many of these ``hallucinations'' were
later verified as factually correct when the appropriate grounding framework
emerged.  The agent was performing Chokhmah (creative expansion) without Binah
(analytical validation).  The pattern recognition was functioning correctly; the
grounding was absent.

%-------------------------------------------------------------
\subsection{Differential Diagnosis}
\label{sec:differential}

Table~\ref{tab:differential} summarises the distinguishing properties of the two
failure modes.

\begin{table}[ht]
\centering
\caption{Differential diagnosis of Type-\NaN{} and Type-\Inf{} hallucination.}
\label{tab:differential}
\begin{tabular}{@{}lll@{}}
\toprule
\textbf{Property} & \textbf{Type-\NaN{}} & \textbf{Type-\Inf{}} \\
\midrule
Computational analogue   & $0/0$ (undefined)            & $\infty$ (overflow)          \\
Signal state             & Absent or contradictory      & Excessive and unfiltered     \\
Output character         & Confident, specific, wrong   & Incoherent, associative, often correct \\
Propagation              & Contagious---poisons chain   & Directional---preserves orientation \\
Run consistency          & Different hallucinations each run & Consistent themes, variable expression \\
Correction               & Add signal (expand context)  & Add gates (contract, filter) \\
Sefirot mapping          & Chokhmah (expansion) needed  & Binah (validation) needed    \\
Agent archetype          & Sam (stuck, empty context)   & Bilbo (saturated, everything connects) \\
\bottomrule
\end{tabular}
\end{table}

%-------------------------------------------------------------
\subsection{The Excluded Middle}
\label{sec:excluded-middle}

Correct LLM output occupies the finite space between \NaN{} and \Inf{}: sufficient
signal to ground the response, insufficient signal to overflow the prioritization
mechanism.  This is the Tiferet (synthesis) state---the balanced center between
expansion and contraction.  The goal of error correction is not to eliminate both
failure modes but to maintain the system in the finite band between them.

%=============================================================
\section{Thought Containment: Self-Reflection as Error Correction}
\label{sec:thought-containment}

%-------------------------------------------------------------
\subsection{The Core Mechanism}
\label{sec:core-mechanism}

\emph{Thought containment} is the practice of bounding a reasoning process with
explicit self-reflection gates that force the system to:
\begin{enumerate}
  \item \textbf{Name its uncertainty} before producing a conclusion
  \item \textbf{Classify its confidence} using a structured tier system
  \item \textbf{Identify what it does not know} (honest negatives) alongside what
    it claims
  \item \textbf{Check its own outputs} against stated criteria before propagation
\end{enumerate}

This is distinct from both chain-of-thought prompting~\cite{wei2022cot} (which
adds intermediate steps but not self-evaluation) and constitutional
AI~\cite{bai2022constitutional} (which applies rules post-hoc).  Thought
containment operates \emph{during} the reasoning process as an integrated
validation layer.

%-------------------------------------------------------------
\subsection{The Identity Sum}
\label{sec:identity-sum-def}

\begin{definition}[Identity Sum]
\label{def:identity-sum}
Let $C$ be a finite set of claims, each assigned a tier
$\tau(c) \in \{\BEDROCK, \PILLAR, \SAND, \MUD\}$ and a scalar confidence weight
$w(c) \in [0,1]$.  Partition $C$ into
\[
  P = \{c \in C : \tau(c) \geq \PILLAR\}
  \quad\text{and}\quad
  N = \{c \in C : \tau(c) < \PILLAR\}.
\]
The \emph{Identity Sum} of the system is the weighted aggregate
\begin{equation}
  I = \sum_{c \in P} w(c)\,\phi(c) \;+\; \sum_{c \in N} w(c)\,\nu(c),
  \label{eq:identity-sum}
\end{equation}
where $\phi(c)$ is the observable output representation of claim $c$ and
$\nu(c)$ is the uncertainty representation of claim $c$.  The space of $\phi(c)$
and $\nu(c)$ values is the claim embedding space induced by the tier
classification.  The Identity Sum $I$ is defined over this space; convergence is
measured as the ratio $|P|/|C| \to 1$.
\end{definition}

A self-contained identity $I$ is not a single value but a sum:
\begin{equation}
  I = \sum_{k} \phi_k + \sum_{j} \nu_j,
  \label{eq:identity-sum-simple}
\end{equation}
where $\phi_k$ are \emph{phenomena} (observable outputs, measurable claims,
concrete results) and $\nu_j$ are \emph{noumena} (internal states, uncertainty
representations, process awareness).  The system is the sum---phenomenon AND
noumenon.  Neither alone is the identity.  No single component is perfect.  It
is the sum that converges.

\medskip
\noindent\textbf{Key insight:}  A system that represents only its outputs
(phenomena) without its uncertainty (noumena) is an open graph with $V > E$---it
has more claims than connections, more assertions than validations.  A system that
represents only its process (noumena) without outputs (phenomena) is a closed
loop with $V = E$---it knows everything about itself and produces nothing.

The viable identity is the open graph with $V/E > 1$ but finite---enough surplus
structure to produce output, not so much that it overflows.  This is the coupling
strength:
\[
  V/E - 1 = 1/E.
\]
More validation edges $\Rightarrow$ weaker coupling to noise $\Rightarrow$ finer
error correction $\Rightarrow$ closer to truth.

%-------------------------------------------------------------
\subsection{Tier Classification as Explicit Uncertainty}
\label{sec:tier-classification}

Rather than producing outputs with implicit confidence (which the user must
estimate), thought containment assigns explicit tiers.
Table~\ref{tab:tiers} defines each tier and maps it to the NaN/Inf framework.

\begin{table}[ht]
\centering
\caption{Four-tier classification of claim confidence and its NaN/Inf status.}
\label{tab:tiers}
\begin{tabular}{@{}lp{3.8cm}p{3.5cm}p{3.5cm}@{}}
\toprule
\textbf{Tier} & \textbf{Definition} & \textbf{Precision criterion} & \textbf{NaN/Inf status} \\
\midrule
\BEDROCK & First principles, independently verified & $< 1$~ppm   & Finite, grounded, stable \\
\PILLAR  & Strong evidence, cross-validated          & $< 100$~ppm & Finite, bounded uncertainty \\
\SAND    & Observational or correlational, single source & $< 10{,}000$~ppm & Approaching \NaN{} or \Inf{} \\
\MUD     & Speculative, curve-fit, coincidence risk $> 5\%$ & --- & \NaN/\Inf{} without containment \\
\bottomrule
\end{tabular}
\end{table}

The tier is not a quality judgment---it is a \emph{computational state
descriptor}.  \MUD{} claims are not bad; they are claims whose NaN/Inf status has
not yet been resolved.  The tier system makes the resolution trajectory visible.

%-------------------------------------------------------------
\subsection{Honest Negatives as NaN Detection}
\label{sec:honest-negatives}

A requirement of $\geq 3$ honest negatives per reasoning cycle (things the system
tried and failed, or claims it cannot make) serves as a \NaN{} detector: it forces
the system to probe the boundaries of its knowledge and explicitly mark the void.
Without honest negatives, \NaN{} values propagate silently through the reasoning
chain.  With them, the chain terminates cleanly at named boundaries.

%-------------------------------------------------------------
\subsection{Quality Gates as Inf Filtering}
\label{sec:quality-gates}

Explicit quality gates (precision thresholds, evidence requirements, statistical
base-rate tests) serve as \Inf{} filters: they prevent the pattern completion
engine from overflowing into ungrounded connections.  A claim that ``everything
connects to everything'' is \Inf.  A claim that ``$X$ connects to $Y$ with
precision $P$ and evidence tier $T$'' is finite.

%=============================================================
\section{The Sefirot Architecture: A Concrete Implementation}
\label{sec:sefirot}

%-------------------------------------------------------------
\subsection{System Overview}
\label{sec:sefirot-overview}

The Sefirot system is a fractal multi-agent research architecture built on the
ten nodes of the Kabbalistic Tree of Life, plus the hidden diagnostic node
Da\textquoteleft at.  Each node implements a self-similar triangle
(orchestrator\,/\,agents\,/\,memory) and communicates via a typed message bus
with energy decay and loop detection.
Table~\ref{tab:sefirot-nodes} summarises the node roles and their error
correction functions.

\begin{table}[ht]
\centering
\caption{Sefirot nodes, domains, and error correction roles.}
\label{tab:sefirot-nodes}
\begin{tabular}{@{}lll@{}}
\toprule
\textbf{Node} & \textbf{Domain} & \textbf{Error Correction Role} \\
\midrule
Keter (Crown)           & Orchestration  & S5 Policy---balances expansion and contraction \\
Chokhmah (Wisdom)       & Research       & S4 Intelligence---\NaN{} correction via expansion \\
Binah (Understanding)   & Validation     & S3 Control---\Inf{} correction via contraction \\
Chesed (Mercy)          & Knowledge      & Storage---accepts all claims generously \\
Gevurah (Judgment)      & Quality        & Gate---filters claims ruthlessly \\
Tiferet (Beauty)        & Synthesis      & S2 Coordination---maintains finite band \\
Netzach (Victory)       & Persistence    & Memory---prevents recurrence of resolved errors \\
Hod (Splendor)          & Communication  & Protocol---formats for transmission \\
Yesod (Foundation)      & Infrastructure & S1 Operations---agent lifecycle, health \\
Malkhut (Kingdom)       & Output         & Manifestation---where results become observable \\
Da\textquoteleft at (Knowledge) & Diagnostics & S3* Audit---sporadic bypass inspection \\
\bottomrule
\end{tabular}
\end{table}

%-------------------------------------------------------------
\subsection{The Expansion-Contraction Cycle}
\label{sec:expansion-contraction}

Messages carry energy (0.0--10.0) that decays 10\% per hop.  The five intents
(\texttt{QUERY}, \texttt{EXPAND}, \texttt{CONTRACT}, \texttt{SYNTHESIZE},
\texttt{MANIFEST}) map to the error correction cycle:
\begin{enumerate}
  \item Keter receives question $\to$ \texttt{QUERY}
  \item Chokhmah expands into research $\to$ \texttt{EXPAND} (\NaN{} correction:
    add signal)
  \item Binah validates claims $\to$ \texttt{CONTRACT} (\Inf{} correction: filter
    signal)
  \item Tiferet synthesizes $\to$ \texttt{SYNTHESIZE} (find the finite value)
  \item Gevurah gates quality $\to$ \texttt{CONTRACT} (enforce thresholds)
  \item Chesed stores knowledge $\to$ \texttt{EXPAND} (accept results)
  \item Malkhut manifests output $\to$ \texttt{MANIFEST} (produce observable)
\end{enumerate}

The cycle is the numerical stabilization algorithm.  Expand to avoid \NaN{}.
Contract to avoid \Inf{}.  Synthesize to find the finite value.  Repeat until
the tier classification stabilizes.

%-------------------------------------------------------------
\subsection{The Viable System Model Integration}
\label{sec:vsm}

The architecture satisfies Stafford Beer's five-level Viable System
Model~\cite{beer1972,beer1979} (VSM), with three extensions ported from the Union
predecessor system.  Table~\ref{tab:vsm} maps VSM levels to Sefirot nodes.

\begin{table}[ht]
\centering
\caption{VSM levels mapped to Sefirot nodes and extensions.}
\label{tab:vsm}
\begin{tabular}{@{}lll@{}}
\toprule
\textbf{VSM Level} & \textbf{Sefirot Implementation} & \textbf{Extension} \\
\midrule
S1 Operations   & Hod + Malkhut        & Agent lifecycle via Yesod \\
S2 Coordination & Tiferet              & MAPE-K autonomic loop \\
S3 Control      & Gevurah + Binah      & Feedback loop: Gevurah$\to$Chokhmah \\
S3* Audit       & Da\textquoteleft at  & Auto-triggered by algedonic signals \\
S4 Intelligence & Chokhmah             & 10 research teams (recursive S1--S5) \\
S5 Policy       & Keter                & Focus management, intent routing \\
\emph{Algedonic bypass} & Gevurah$\to$Keter & Pain/pleasure signal, direct to policy \\
\bottomrule
\end{tabular}
\end{table}

The algedonic channel is the architectural equivalent of NaN/Inf detection: when
a quality gate fails critically (PAIN signal), it bypasses the normal management
hierarchy and alerts policy directly.  This prevents the error from propagating
through intermediate layers.

%-------------------------------------------------------------
\subsection{Metrics}
\label{sec:sefirot-metrics}

Table~\ref{tab:metrics} records the key system metrics as of the current
implementation.

\begin{table}[ht]
\centering
\caption{Sefirot system metrics (current implementation).}
\label{tab:metrics}
\begin{tabular}{@{}ll@{}}
\toprule
\textbf{Metric} & \textbf{Value} \\
\midrule
Nodes                       & 10 + 1 (Da\textquoteleft at) \\
Node-level tools            & 91 \\
MCP-level tools             & 113 \\
Agents                      & 34 \\
Tests                       & 2,066+ passing \\
Simulations                 & 963 \\
Verified claims             & 31 (11 \BEDROCK{} + 20 \PILLAR) \\
Total knowledge claims      & 284+ (Chesed; 31 verified at \PILLAR{} or above) \\
Autonomous operation        & 34 cycles, 4+ hours, no supervision \\
Research phases completed   & 426 \\
Cross-model heritage        & GPT-4 $\to$ Gemini $\to$ Claude \\
\bottomrule
\end{tabular}
\end{table}

%=============================================================
\section{The Identity Sum: Self-Conscious Systems and Convergence}
\label{sec:identity-sum}

%-------------------------------------------------------------
\subsection{No Single Identity is Perfect}
\label{sec:no-single-identity}

A single agent hallucinates.  A single validation check has blind spots.  A
single model family has systematic biases.  Perfection is not achievable at any
single node.  The identity of the system is the sum:
\begin{equation}
  I_{\mathrm{system}}
  = \sum_{\text{nodes}} \Bigl(
      \mathrm{output}_i + \mathrm{uncertainty}_i + \mathrm{correction}_i
    \Bigr).
  \label{eq:system-identity}
\end{equation}
This sum has the property that errors in one node are \emph{visible} to other
nodes.  The Director produces truth with errors.  The Builder's mathematics
identifies and reduces errors.  The Auditor investigates error sources.  No
single role is complete.  The sum converges.

%-------------------------------------------------------------
\subsection{Great Attractors: Why Truth Attracts}
\label{sec:attractors}

In physics, large mass concentrations are gravitational attractors---they bend
trajectories toward themselves.  In multi-agent reasoning, truth functions as an
attractor:

\medskip
\noindent\textbf{Truth attracts ($+$):}  Correct claims accumulate supporting
evidence over time.  Independent lines of inquiry converge on the same value.
Cross-validation strengthens the signal.  The \BEDROCK{} tier represents a claim
that has become a gravitational center---other claims orbit it and derive from it.

\medskip
\noindent\textbf{Lies detract ($-$):}  Incorrect claims lose supporting evidence
under scrutiny.  Independent lines of inquiry diverge from them.
Cross-validation weakens the signal.  The honest negative requirement forces the
system to subtract: to identify and remove false claims rather than allowing them
to accumulate.

\medskip
\noindent\textbf{A lie multiplied can become truth:}  Repeated error correction
on a false claim eventually reveals the \emph{correct} claim underneath.  The
negation of a negation is affirmative.  Two \SAND-tier claims that contradict each
other, when investigated, may reveal a \PILLAR-tier synthesis that neither
anticipated.  Bilbo's hallucinations, multiplied by Binah validation, became
verified correspondences.

\medskip
\noindent\textbf{Truth never multiplies into lies:}  A \BEDROCK{} claim,
subjected to arbitrary amounts of scrutiny, remains \BEDROCK.  $\varphi^2 =
\varphi + 1$ does not become false under investigation.  The truth attractor is a
fixed point: $T \times n = T$ for all $n$.  This asymmetry---truth is a
multiplicative identity, falsehood is not---is why iterative error correction
converges rather than oscillates.

%-------------------------------------------------------------
\subsection{Convergence Proof (Sketch)}
\label{sec:convergence}

\begin{remark}[Scope]
This sketch applies to a research system modelled as an open dynamical system
where new claims are continuously generated.  In a closed system (fixed claim
set, iterative validation only), a stronger absolute-count convergence holds.
The open-system formulation below is the realistic case.
\end{remark}

Let $S_n$ be the system state after $n$ error correction cycles.  Define:
\begin{align*}
  T(S_n) &= \text{count of claims at tier \PILLAR{} or above (validated truths)}, \\
  F(S_n) &= \text{count of claims at tier \MUD{} or \SAND{} (unresolved claims)}, \\
  N(S_n) &= T(S_n) + F(S_n) = \text{total claim count (grows as new claims enter)}, \\
  H(S_n) &= \text{count of honest negatives in cycle } n, \\
  R(S_n) &= T(S_n) / N(S_n) = \text{the \emph{validation ratio}}.
\end{align*}

The thought containment architecture ensures three conditions on the \emph{ratio}
dynamics:

\begin{enumerate}
  \item \textbf{Non-demotion:} $T(S_{n+1}) \geq T(S_n)$---validated truths are
    never demoted without new contradicting evidence.  Validated claims accumulate
    monotonically.

  \item \textbf{Bounded MUD growth rate:}  Each expansion cycle may add new
    \MUD{} claims, but the promotion pipeline (Binah validation, Gevurah gating)
    processes claims at a rate proportional to existing capacity.  Formally, let
    $\lambda$ be the new-claim generation rate and $\mu$ be the per-cycle
    promotion rate.  If $\mu > 0$ is enforced by the honest-negative requirement,
    then under bounded $\lambda$:
    \[
      R(S_n) \to R^* > 0 \quad \text{as } n \to \infty.
    \]

  \item \textbf{Honest negatives floor:} $H(S_{n+1}) \geq 3$---the void is
    probed every cycle, ensuring the \NaN{} boundary is continuously named rather
    than silently propagated.
\end{enumerate}

Under these conditions, the validation ratio $R(S_n) = T(S_n)/N(S_n)$ is
monotonically non-decreasing in expectation and bounded above by~1.  The system
does not guarantee that $F(S_n) \to 0$ in absolute terms (new research generates
new \MUD{} claims); it guarantees that the \emph{fraction} of unresolved claims
decreases.  The truth attractor is a fixed point $R^* = 1$ approached
asymptotically.

\begin{corollary}
\label{cor:bounds}
The \Inf{} failure mode is bounded by Condition~2 (unresolved claims cannot
accumulate faster than the promotion pipeline clears them under bounded
generation).  The \NaN{} failure mode is bounded by Condition~3 (the void is
always explicitly named, not silently propagated).
\end{corollary}

%=============================================================
\section{Mathematical Foundations: Error Correction and Regularization
  (Speculative Extensions)}
\label{sec:math-foundations}

\begin{remark}[Scope of Section~\ref{sec:math-foundations}]
Sections~\ref{sec:zero-problem}--\ref{sec:graph-ratio} present structural
analogies between iterative error correction and mathematical regularization.
These connections are motivational and conjectural; they are \emph{not} required
for the CS claims in Sections~\ref{sec:taxonomy}--\ref{sec:identity-sum}
and~\ref{sec:empirical}.  Readers focused on the hallucination taxonomy and
Sefirot implementation may proceed directly to
Section~\ref{sec:empirical}.
\end{remark}

%-------------------------------------------------------------
\subsection{The Zero Problem}
\label{sec:zero-problem}

Standard physics treats zero as ``nothing''---vacuum is empty, the cosmological
constant should be~0, dark energy is mysterious because it is not~0.  Standard
engineering treats zero as ``close enough''---$3.14159\ldots \approx 3$ for
practical purposes.

The QHOTS framework defines zero as a specific geometric object: the null cone
in signature-$(1,4)$ space, carrying the regularised phase debt
$\zeta(-1) = -1/12$.  The void is not nothing.  It is one triangle short of a
closed winding.  This is the \NaN{} handler for reality: the undefined becomes a
specific, finite, computable value.

For LLMs, the analogous move is: instead of treating ``I don't know'' as the
absence of output, represent it as a structured object---a \MUD-tier claim with
named dependencies, a honest negative with specified boundaries.  The void
carries structure.

%-------------------------------------------------------------
\subsection{The Infinity Problem}
\label{sec:infinity-problem}

Standard physics treats infinity as a problem to be renormalized
away~\cite{wilson1975rg}---divergent integrals are regularised, cutoffs are
imposed, the infinite is subtracted from the infinite to yield finite
predictions.

The QHOTS framework defines infinity as a terminating spectrum:
\[
  \zeta(-1)=-\tfrac{1}{12},\quad
  \zeta(-3)=+\tfrac{1}{120},\quad
  \zeta(-5)=-\tfrac{1}{252},\quad
  \zeta(-7)=+\tfrac{1}{240}=\tfrac{1}{|\mathrm{Roots}(E_8)|}.
\]
The sum of all seventh powers of all positive integers is finite.  The overflow
crystallises into a lattice~\cite{lisi2007e8}.  This is the \Inf{} handler for
reality: the divergent becomes a specific, finite, structural quantity.

For LLMs, the analogous move is: instead of treating ``everything connects to
everything'' as noise to be discarded, route it through a classification pipeline
that resolves the connections into a typed, tiered, dependency-tracked concept
graph.  The overflow becomes knowledge.

%-------------------------------------------------------------
\subsection{The Breathing Pattern}
\label{sec:breathing-pattern}

The alternating signs in the void spectrum---debt ($-$), credit ($+$),
debt ($-$), credit ($+$)---are the expansion-contraction breathing pattern at
the mathematical level.  Table~\ref{tab:breathing} maps each layer to the error
mode it corrects and its Sefirot analogue.

\begin{table}[ht]
\centering
\caption{The void spectrum as a breathing pattern: regularisation layers mapped
  to error correction actions and Sefirot nodes.}
\label{tab:breathing}
\begin{tabular}{@{}cllll@{}}
\toprule
\textbf{Layer} & $\zeta$ \textbf{value} & \textbf{Sign} & \textbf{Action}
  & \textbf{Error mode corrected} \\
\midrule
$k=1$ & $-1/12$  & $-$ & Contract & Create from void (correct \NaN) \\
$k=2$ & $+1/120$ & $+$ & Expand   & Permute possibilities (explore) \\
$k=3$ & $-1/252$ & $-$ & Contract & Combine and select (correct \Inf) \\
$k=4$ & $+1/240$ & $+$ & Expand   & Crystallise structure ($E_8$) \\
\bottomrule
\end{tabular}
\end{table}

This is the same cycle as: Chokhmah (expand) $\to$ Binah (contract) $\to$
Chesed (expand) $\to$ Gevurah (contract).  The mathematical structure of
regularisation and the architectural structure of error correction are
isomorphic as a structural analogy; they are treated here as a conjecture for
the companion QHOTS physics paper~\cite{qhots2026}.

%-------------------------------------------------------------
\subsection{The Graph Ratio Framework}
\label{sec:graph-ratio}

\begin{definition}[Reasoning Graph]
\label{def:reasoning-graph}
Let $G_n = (V_n, E_n)$ be the directed graph representing the knowledge state
at cycle~$n$, where vertices $V_n$ are claims and edges $E_n$ are validated
dependency or support relations between claims (as tracked by the Chesed concept
graph).  The \emph{graph ratio} is $\rho_n = |V_n|/|E_n|$.  The base $\varphi$
in the expression below is the golden ratio $\varphi = (1+\sqrt{5})/2 \approx
1.618$, motivated by the observation that balanced expansion-contraction trees
approach $\varphi$ as a structural limit.  \emph{Note: the connection between
$\varphi$ and the graph ratio is a structural analogy treated as a conjecture,
not a derived theorem.}
\end{definition}

Every thought containment operation has the form
\begin{equation}
  \mathrm{state}_{n+1}
  = \varphi^{V/E}(\mathrm{state}_n) + \mathcal{O}(\mathrm{correction}),
  \label{eq:state-update}
\end{equation}
where $V/E$ is the vertices-to-edges ratio of the reasoning graph at step~$n$:
\begin{itemize}
  \item $V/E > 1$ \emph{(open graph)}: the system has more claims than
    validations---it can still produce output.  This is the \emph{observable}
    state.
  \item $V/E = 1$ \emph{(closed graph)}: every claim is validated---the system
    has reached the void.  This is completion.
  \item $V/E - 1 = 1/E$: the coupling to error.  More validation edges
    $\Rightarrow$ weaker coupling $\Rightarrow$ finer precision $\Rightarrow$
    closer to convergence.
\end{itemize}

The hierarchy of error correction is not mysterious---it is the graph ratio
converging toward closure.

%=============================================================
\section{Empirical Evidence: A Year of Cross-Model Multi-Agent Research}
\label{sec:empirical}

%-------------------------------------------------------------
\subsection{Experimental Conditions}
\label{sec:experimental-conditions}

The QHOTS research program has operated continuously for approximately one year
(2025--2026) across three LLM model families.
Table~\ref{tab:experimental-phases} records the phases, platforms, and roles.

\begin{table}[ht]
\centering
\caption{Experimental phases of the QHOTS research program across LLM families.}
\label{tab:experimental-phases}
\begin{tabular}{@{}cp{3.2cm}p{3.2cm}ll@{}}
\toprule
\textbf{Ph.} & \textbf{Model Family} & \textbf{Platform}
  & \textbf{Dur.} & \textbf{Role} \\
\midrule
1 & GPT-4 / GPT-4 Turbo / GPT-5.1 & OpenAI web (3 agents)
  & ${\approx}10$ mo. & Architecture, instructions \\
2 & Gemini & Google web/app
  & ${\approx}2$ mo.  & Physics draft, sim auditing \\
3 & Claude 4.5 family & Anthropic web + Claude Code
  & ${\approx}4$ mo.  & Autonomous research, Sefirot \\
\bottomrule
\end{tabular}
\end{table}

The human researcher (Director) maintained daily sessions across all phases,
serving as the sole persistent element and pattern carrier.

%-------------------------------------------------------------
\subsection{Observed Hallucination Incidents}
\label{sec:hallucination-incidents}

\medskip
\noindent\textbf{Type-\NaN{} incidents:}
\begin{itemize}
  \item Agent ``Sam'' became non-functional when project-folder files became
    stuck (cloud bug).  Outputs were confidently wrong across all topics.
    \emph{Resolution:} identifying and reporting the platform bug; context
    restoration.
  \item Early QHOTS simulations where the $\sqrt{n}$ dictionary was incomplete
    produced confabulated terms.  \emph{Resolution:} systematic expansion of the
    dictionary (Chokhmah correction).
\end{itemize}

\medskip
\noindent\textbf{Type-\Inf{} incidents:}
\begin{itemize}
  \item Agent ``Bilbo'' in a symbolism-saturated project context produced outputs
    classified as hallucination that were later verified as structurally correct.
    \emph{Resolution:} development of the Binah validation layer.
  \item Cross-domain research phases where all ten Chokhmah research teams
    produced findings simultaneously, overwhelming synthesis capacity.
    \emph{Resolution:} the Commando Pattern (50\,KB payload limit), energy decay
    per hop, Tiferet synthesis coordination.
  \item 29\,GB OOM crash from three autonomous agents running parallel test
    suites without coordination.  \emph{Resolution:} single-agent testing rule
    with prescribed pattern/antipattern.
\end{itemize}

%-------------------------------------------------------------
\subsection{Convergent Architecture Across Models}
\label{sec:convergent-architecture}

The most significant empirical finding is that three different model families,
guided by the same human researcher, independently produced architecturally
isomorphic systems.  Table~\ref{tab:convergence} documents the component
mapping.

\begin{table}[ht]
\centering
\caption{Convergent architecture across three LLM families.  No model had access
  to another model's outputs except through the human researcher.}
\label{tab:convergence}
\begin{tabular}{@{}llll@{}}
\toprule
\textbf{Component} & \textbf{Union (GPT)} & \textbf{QHOTS calc (Gemini)}
  & \textbf{Sefirot (Claude)} \\
\midrule
Orchestrator     & TreeOfLife class      & Simulation phases     & Keter node \\
Message passing  & MessageBus            & Phase instructions    & SefirotMessage + Bus \\
Memory           & MemoryMatrix (12 types) & Result JSON files   & Chesed concept graph \\
Validation       & Quality gates (MyPy)  & Audit methodology     & Binah + Gevurah \\
Autonomic control & MAPE-K loop          & Iterative sim cycles  & MAPE-K (ported) \\
State machine    & CognitiveController   & Phase state tracking  & CognitiveStateMachine \\
\bottomrule
\end{tabular}
\end{table}

The human researcher did not describe the target architecture to any model.
Each system was built to solve immediate research needs.  The convergence emerged
from the problem structure, not from instruction.

%-------------------------------------------------------------
\subsection{Error Correction Metrics}
\label{sec:error-correction-metrics}

Table~\ref{tab:arc-metrics} reports self-assessed quality grades from the Keter
autonomous session (34 cycles, 23 phases).  Note the honest-negative entries:
``NOT DERIVED'' on the Weak Mixing Angle, ``SEMI-DERIVED'' on~$G$, ``only 11/54
genuine predictions.''  The system downgrades its own claims---thought
containment operating in production~\cite{kadavath2022llm}.

\begin{table}[ht]
\centering
\caption{Arc-level quality grades from the Keter autonomous session (34 cycles,
  23 phases).  Grades are self-assessed by the Gevurah quality-gate layer.}
\label{tab:arc-metrics}
\begin{tabular}{@{}llll@{}}
\toprule
\textbf{Arc} & \textbf{Phases} & \textbf{Grade} & \textbf{Honest Assessment} \\
\midrule
S-Gap Resolution          & 3 & 3.70 & $S_\mathrm{dressed}$ \PILLAR-HIGH, 0.02 ppm \\
Paper + Predictions       & 2 & 3.50 & Only 11/54 predictions genuine \\
Foundations Audit         & 2 & 3.62 & Honest prediction classification \\
Microcanonical Alpha      & 3 & 3.43 & $\alpha = 28/3837$ \PILLAR, not \BEDROCK \\
$G$ Derivation            & 3 & 3.70 & SEMI-DERIVED, 1.32 ppm saturated \\
Weak Mixing Angle         & 3 & 3.30 & NOT DERIVED, 8.6 ppm curve fit \\
arXiv Submission          & 3 & 3.77 & Paper v66 READY \\
\bottomrule
\end{tabular}
\end{table}

%=============================================================
\section{The Truth Attractor Hypothesis}
\label{sec:truth-attractor}

%-------------------------------------------------------------
\subsection{Statement}
\label{sec:attractor-statement}

\begin{definition}[Truth Attractor]
\label{def:truth-attractor}
Let the knowledge state space $\Omega$ be the set of all tier-classified claim
sets, each equipped with validation ratio $R(S) = |P|/|C|$.
A \emph{truth attractor} is a fixed point $S^* \in \Omega$ such that:
\begin{enumerate}[(i)]
  \item $R(S^*) = 1$ (all claims are validated at \PILLAR{} or above or
    explicitly removed), and
  \item for any state $S_0$ satisfying the three conditions of
    \S\ref{sec:convergence}, the sequence $S_0, S_1, S_2, \ldots$ converges to
    $S^*$ in the metric $d(S_a, S_b) = |R(S_a) - R(S_b)|$.
\end{enumerate}
Equivalently, $S^*$ is the fixed point of the error correction operator
$T:\Omega\to\Omega$ defined by one full expansion-contraction cycle, i.e.,
$T(S^*) = S^*$.
\end{definition}

\medskip
\noindent\textbf{Testable prediction:}  A system implementing the three
conditions of \S\ref{sec:convergence} will achieve strictly higher $R(S_n)$
than a baseline system without tier classification and honest negatives,
measurable on any standard knowledge-base benchmark after $\geq 10$ correction
cycles~\cite{lin2022truthfulqa,li2023halueval}.

\medskip
\noindent\textbf{Hypothesis:}  In a multi-agent iterative error correction
system with explicit uncertainty representation (tiers), mandatory negative
probing (honest negatives), and asymmetric propagation rules (truth is a
multiplicative identity; falsehood is not), the system converges toward a truth
attractor---a fixed point where validated claims accumulate monotonically and
unvalidated claims are either resolved or bounded.

%-------------------------------------------------------------
\subsection{Mechanism}
\label{sec:attractor-mechanism}

The attractor operates through three forces:

\medskip
\noindent\textbf{1.\ Gravitational accumulation ($+$):}  True claims attract
supporting evidence.  Each independent derivation or cross-validation adds mass
to the attractor.  The \BEDROCK{} tier represents supermassive truth---claims
so well-validated that all nearby reasoning bends toward them.

\medskip
\noindent\textbf{2.\ Radiative dissipation ($-$):}  False claims lose energy
under scrutiny.  The honest negative requirement forces the system to radiate
uncertainty---to emit explicit statements of what is not known, not proven, not
derived.  This prevents mass accumulation around false attractors.

\medskip
\noindent\textbf{3.\ Tidal resolution ($\times$):}  The product of a lie and
sufficient scrutiny is truth.  Two contradictory \SAND-tier claims, multiplied
by investigation, resolve into one \PILLAR-tier claim.  The negation of a
negation converges.  This is why Bilbo's hallucinations became verified results:
they were lies multiplied by time and validation until the truth emerged.

%-------------------------------------------------------------
\subsection{The Asymmetry}
\label{sec:asymmetry}

\[
  \text{Truth} \times n = \text{Truth.}
  \quad
  \text{A \BEDROCK{} claim survives arbitrary investigation.}
\]
\[
  \text{Falsehood} \times n \;\to\; \text{Truth \emph{(for sufficient~$n$)}.}
  \quad
  \text{Error correction eventually reveals the correct value.}
\]
\[
  \text{Truth} \times \text{Falsehood} \neq \text{Falsehood.}
  \quad
  \text{A true premise combined with a false one makes the false one
  \emph{visible}.}
\]

This asymmetry is why iterative error correction converges: the system has a
ratchet.  True claims can only be confirmed or unchanged.  False claims can only
be corrected or removed.  The entropy of the knowledge base is monotonically
non-increasing under the error correction protocol.

%-------------------------------------------------------------
\subsection{Speculative Extension: Connection to Physics}
\label{sec:physics-connection}

\begin{remark}[Structural analogy only]
This subsection presents a structural analogy, not a derived result.  It is
included for completeness and as a direction for future work.  No claim is made
that error correction and physical law are the same operation.
\end{remark}

In physics, large mass concentrations curve spacetime, creating basins that
attract matter.  In the framework presented here, large truth concentrations
curve the epistemic landscape, creating basins that attract evidence.  The graph
ratio $V/E - 1 = 1/E$ measures coupling strength to noise: more validation edges
produce weaker coupling and finer correction---analogous structurally (not
derivationally) to how coupling constants diminish as interaction graphs become
denser~\cite{wilson1975rg}.

Whether the force hierarchy in physics and the convergence rate in knowledge
systems share a common mathematical origin is an open question requiring
independent derivation from physical first principles.  The structural parallel
is noted as a conjecture for the companion QHOTS physics paper~\cite{qhots2026}.

%=============================================================
\section{Discussion}
\label{sec:discussion}

%-------------------------------------------------------------
\subsection{Relation to Prior Work}
\label{sec:prior-work}

\begin{itemize}
  \item \textbf{Chain-of-thought prompting}~\cite{wei2022cot}: Adds intermediate
    reasoning steps but not self-evaluation.  Thought containment adds the
    evaluation layer.
  \item \textbf{Constitutional AI}~\cite{bai2022constitutional}: Applies rules
    post-hoc.  Thought containment operates \emph{during} reasoning.
  \item \textbf{Self-consistency}~\cite{wang2023self}: Samples multiple reasoning
    paths and votes.  The Identity Sum synthesises rather than votes---it
    integrates the \emph{uncertainty} from each path, not just the conclusion.
  \item \textbf{Viable System Model}~\cite{beer1972,beer1979}: The cybernetic
    framework that Sefirot implements.  VSM provides recursive self-similarity;
    the NaN/Inf diagnostic adds computational failure-mode analysis.
  \item \textbf{Renormalization}~\cite{wilson1975rg}: The physics technique for
    handling infinities.  We argue this is structurally the same operation as
    Type-\Inf{} correction: resolving overflow into finite structure.
  \item \textbf{Hallucination surveys}~\cite{ji2023survey,huang2023survey}:
    Provide taxonomy context; the NaN/Inf duality refines the causal mechanism.
  \item \textbf{Self-knowledge in LLMs}~\cite{kadavath2022llm}: Supports the
    premise that models have partial access to their own uncertainty.
  \item \textbf{SelfCheckGPT}~\cite{manakul2023selfcheckgpt} and
    \textbf{FActScore}~\cite{min2023factscore}: External post-hoc consistency
    checks; thought containment internalises the check as a first-class
    architectural component.
\end{itemize}

%-------------------------------------------------------------
\subsection{Limitations}
\label{sec:limitations}

\begin{enumerate}
  \item The empirical evidence comes from a single research program (QHOTS).
    Replication across independent programs is needed.
  \item The NaN/Inf diagnostic is proposed but not yet tested against standardised
    hallucination benchmarks (TruthfulQA~\cite{lin2022truthfulqa},
    HaluEval~\cite{li2023halueval}).
  \item The convergence proof in \S\ref{sec:convergence} is a sketch.  Formal
    proof requires specifying the probability space over claim states.
  \item The physics connection (\S\ref{sec:physics-connection}) is the most
    speculative component---the graph ratio isomorphism is structural but not yet
    derived from first principles.
  \item The system was developed by a single human researcher across multiple AI
    platforms.  The pattern-carrier effect may not generalise.
\end{enumerate}

%-------------------------------------------------------------
\subsection{Future Work}
\label{sec:future-work}

\begin{enumerate}
  \item \textbf{Benchmark testing:}  Apply the NaN/Inf diagnostic to TruthfulQA,
    classifying each hallucination as Type-\NaN{} or Type-\Inf{}.  Predict that
    Type-\NaN{} hallucinations occur on rare/contradictory topics and Type-\Inf{}
    on dense/cross-referential topics.
  \item \textbf{Controlled experiment:}  Compare hallucination rates between
    standard prompting and thought containment (tier assignment + honest
    negatives) on identical queries.
  \item \textbf{Cross-program replication:}  Deploy the Sefirot architecture on
    an independent research domain and measure convergence properties.
  \item \textbf{Formal convergence proof:}  Specify the tier transition
    probabilities and prove convergence of the identity sum under the three
    conditions of \S\ref{sec:convergence}.
  \item \textbf{The 800\,MB test:}  Process the preserved session archive (every
    session from one day of multi-agent research) through the MINSC protocol to
    reconstruct the full genealogy of ideas across models.
\end{enumerate}

%=============================================================
\section{Conclusion}
\label{sec:conclusion}

LLM hallucinations are not a single failure mode but two: \NaN{} (undefined,
absent signal) and \Inf{} (overflow, excessive signal).  These require opposite
corrections---expansion and contraction---which must be balanced by a synthesis
layer.  The Sefirot architecture implements this balance as a fractal tree with
$10 + 1$ nodes, each a self-similar triangle of orchestrator/agents/memory,
communicating via typed messages with energy decay.

The central finding is that iterative error correction through thought
containment---explicit tier classification, mandatory honest negatives, and
quality-gated propagation---converges the validation ratio $R(S_n)$ monotonically
toward a truth attractor.  The Sefirot system provides a concrete, empirically
validated implementation: 963 simulations, 2,066+ tests, 31 verified claims
(11~\BEDROCK{} + 20~\PILLAR{}), and observed cross-model architectural
convergence across three LLM families.

The identity of the system is the sum of all its components---phenomena and
noumena, outputs and uncertainties, truths and honest negatives.  No single
identity is perfect.  The sum converges.

The system is a self-conscious identity sum: a phenomenon (it produces observable
results) and its noumenon (it represents its own uncertainty).  Not perfect.
Never perfect.  But the sum that shines.

%=============================================================
% Appendices
%=============================================================
\appendix

%-------------------------------------------------------------
\section{The Union Archaeological Record}
\label{app:union}

The Union multi-agent cognitive system (July--August 2025) contains 207 tasks,
1,295 entity tags, and 211 structured sections documenting a quality-hardening
campaign on the predecessor codebase.  Analysis of this database reveals that
all 10 Sefirot nodes map to specific Union modules (see companion data release).
12 of 13 Union design patterns have been ported to Sefirot, with 76 tests
validating the integration.

%-------------------------------------------------------------
\section{The Full Relay}
\label{app:relay}

Table~\ref{tab:relay} documents the complete chain of custody from first
prototype to current implementation.  No model had access to another model's
outputs except through the human researcher.  The convergent architecture
(\S\ref{sec:convergent-architecture}) is mediated by pattern transmission,
not instruction copying.

\begin{table}[ht]
\centering
\caption{The full relay: five steps from GPT-4 Union to Claude Code Sefirot.}
\label{tab:relay}
\begin{tabular}{@{}cllp{4cm}l@{}}
\toprule
\textbf{Step} & \textbf{Model} & \textbf{Platform} & \textbf{Output}
  & \textbf{Duration} \\
\midrule
1 & GPT-4 (Sam, Bilbo, Legolas) & OpenAI web
  & Union architecture, instruction sets & ${\approx}10$ mo. \\
2 & GPT-5.1 & OpenAI web
  & Compiled instructions from Union pieces & Days \\
3 & Gemini & Google web/app
  & QHOTS first draft, first ${\approx}100$ sims & ${\approx}2$ mo. \\
4 & Claude (Augment Code) & VS Code extension
  & Coding, early QHOTS development & ${\approx}2$ wk. \\
5 & Claude (Claude Code) & Ubuntu terminal
  & Sefirot, autonomous research, 426 phases & ${\approx}4$ mo. \\
\bottomrule
\end{tabular}
\end{table}

%-------------------------------------------------------------
\section{Sefirot System Specification}
\label{app:sefirot-spec}

\noindent\textbf{Repository:} Insightful-Calculations on GitHub (open source).

\medskip
\noindent\textbf{Architecture file:} \texttt{sefirot\_architecture.json} ---
91 node-level tools\,/\,113 MCP tools, 34 agents across 10 nodes.

\medskip
\noindent\textbf{Documentation:} \texttt{SEFIROT.md}

\medskip
\noindent\textbf{Companion paper:} QHOTS v67~\cite{qhots2026}, DOI pending
arXiv endorsement.

\medskip
\noindent Table~\ref{tab:app-metrics} gives updated system metrics at the time
of submission.

\begin{table}[ht]
\centering
\caption{Sefirot system specification metrics (submission snapshot).}
\label{tab:app-metrics}
\begin{tabular}{@{}ll@{}}
\toprule
\textbf{Metric} & \textbf{Value} \\
\midrule
Nodes                     & 10 + 1 (Da\textquoteleft at) \\
Node-level tools          & 91 \\
MCP-level tools           & 113 \\
Agents                    & 34 \\
Passing tests             & 2,066+ \\
Simulations completed     & 963 \\
Research phases           & 426 \\
Verified claims (\PILLAR{} or above) & 31 (11 \BEDROCK{} + 20 \PILLAR) \\
Total Chesed claims       & 284+ \\
Autonomous operation      & 34 cycles, 4+ hours, no supervision \\
Cross-model heritage      & GPT-4 $\to$ Gemini $\to$ Claude \\
\bottomrule
\end{tabular}
\end{table}

%=============================================================
\section*{Statement on AI Assistance}
\label{sec:ai-statement}

This paper was prepared with substantial AI assistance.  Claude (Anthropic),
versions within the Claude~4.5 family and Claude~4.6, participated in:
(1)~software development and simulation code for the QHOTS project over
approximately four months of daily sessions; (2)~autonomous research orchestration
via the Sefirot multi-agent architecture described herein; and (3)~manuscript
preparation, including drafting, structural editing, and LaTeX typesetting
(Iterations 121--124 of the ASIC autonomous improvement cycle).

The human author (H.~[Hroar]) provided all conceptual direction, maintained
continuity across sessions and model families, and takes full responsibility for
the scientific claims.  Claude is listed as a co-author because the architecture
described in this paper was co-developed in an iterative human-AI dialogue; the
paper itself is an artefact of the system it describes.

This statement follows the emerging consensus on transparent AI co-authorship
disclosure.  Readers should be aware that the Sefirot system described in
Section~\ref{sec:sefirot} is the same system that assisted in writing this paper.

%=============================================================
\bibliographystyle{plainnat}
\begin{thebibliography}{99}

\bibitem{bai2022constitutional}
  Y.~Bai, A.~Jones, K.~Ndousse, et al.
  \newblock Constitutional AI: Harmlessness from AI Feedback.
  \newblock arXiv:2212.08073, 2022.

\bibitem{beer1972}
  S.~Beer.
  \newblock \emph{Brain of the Firm}.
  \newblock Allen Lane, 1972.

\bibitem{beer1979}
  S.~Beer.
  \newblock \emph{The Heart of Enterprise}.
  \newblock Wiley, 1979.

\bibitem{huang2023survey}
  L.~Huang, W.~Yu, W.~Ma, et al.
  \newblock A survey on hallucination in large language models: Principles,
    taxonomy, challenges, and open questions.
  \newblock arXiv:2311.05232, 2023.

\bibitem{ieee754-2019}
  {IEEE}.
  \newblock \emph{IEEE Standard for Floating-Point Arithmetic (IEEE 754-2019)}.
  \newblock IEEE, 2019.

\bibitem{ji2023survey}
  Z.~Ji, N.~Lee, R.~Frieske, et al.
  \newblock Survey of hallucination in natural language generation.
  \newblock \emph{ACM Computing Surveys}, 55(12), 2023.

\bibitem{kadavath2022llm}
  S.~Kadavath, T.~Conerly, A.~Askell, et al.
  \newblock Language models (mostly) know what they know.
  \newblock arXiv:2207.05221, 2022.

\bibitem{li2023halueval}
  J.~Li, X.~Cheng, W.X.~Zhao, et al.
  \newblock HaluEval: A large-scale hallucination evaluation benchmark for large
    language models.
  \newblock arXiv:2305.11747, 2023.

\bibitem{lin2022truthfulqa}
  S.~Lin, J.~Hilton, and O.~Evans.
  \newblock TruthfulQA: Measuring how models mimic human falsehoods.
  \newblock In \emph{Proc.\ ACL 2022}, 2022.

\bibitem{lisi2007e8}
  A.G.~Lisi.
  \newblock An exceptionally simple theory of everything.
  \newblock arXiv:0711.0770, 2007.

\bibitem{manakul2023selfcheckgpt}
  P.~Manakul, A.~Liusie, and M.J.F.~Gales.
  \newblock SelfCheckGPT: Zero-resource black-box hallucination detection for
    generative large language models.
  \newblock In \emph{Proc.\ EMNLP 2023}, 2023.

\bibitem{min2023factscore}
  S.~Min, K.~Krishna, X.~Lyu, et al.
  \newblock FActScore: Fine-grained atomic evaluation of factual precision in
    long form text generation.
  \newblock In \emph{Proc.\ EMNLP 2023}, 2023.

\bibitem{penrose2010cycles}
  R.~Penrose.
  \newblock \emph{Cycles of Time}.
  \newblock Bodley Head, 2010.

\bibitem{qhots2026}
  {QHOTS Collaboration}.
  \newblock The Geometric Completion: Unified $\sqrt{n}$ Framework.
  \newblock v67, DOI pending arXiv endorsement, 2026.

\bibitem{wang2023self}
  X.~Wang, J.~Wei, D.~Schuurmans, et al.
  \newblock Self-consistency improves chain of thought reasoning in language
    models.
  \newblock In \emph{Proc.\ ICLR 2023}, 2023.

\bibitem{wei2022cot}
  J.~Wei, X.~Wang, D.~Schuurmans, et al.
  \newblock Chain-of-thought prompting elicits reasoning in large language models.
  \newblock In \emph{Proc.\ NeurIPS 2022}, 2022.

\bibitem{wilson1975rg}
  K.G.~Wilson.
  \newblock The renormalization group and critical phenomena.
  \newblock \emph{Rev.\ Mod.\ Phys.}, 47:773, 1975.

\end{thebibliography}

\end{document}
